% Part: normal-modal-logic
% Chapter: filtrations
% Section: S5-decidable
%READERNOTE{SOL6-A172: निर्णेयतेच्या प्रस्तावनेत दिलेली सांत आणि प्रभावी रूपबंध-यादी स्पष्ट केली. S5 च्या प्रतिमान-शोधात प्रत्येक आकारासाठी निश्चित जगांची नावे, वैश्विक संबंध आणि लक्ष्य सूत्रातील चलांचीच सांत मूल्यांकन-यादी वापरली जाते; उर्वरित चलांचे मूल्यांकन रिक्त ठेवता येते.}

\documentclass[../../../include/open-logic-section]{subfiles}

\begin{document}

\olfileid{nml}{fil}{dec}

\olsection{\Log{S5} निर्णेय आहे}
दिलेल्या सांत रूपबंध-यादीने प्रभावीपणे परिभाषित मोडल तर्कशास्त्राच्या
प्रणाली \emph{निर्णेय} आहेत हे दाखवण्याचा
सांत प्रतिमान गुणधर्म हा सोपा मार्ग देतो.
म्हणजे, !!a{formula} त्या प्रणालीत
!!{derivable} आहे की नाही हे ठरवणारी
संगणनक्षम पद्धत मिळते.

\begin{thm}
  \Log{S5} निर्णेय आहे.
\end{thm}

\begin{proof}
  $!A$ दिले आहे असे समजा आणि~$!A$ मध्ये येणारी
  !!{propositional variable}s ही $p_1$, \dots, $p_k$
  यांपैकीच आहेत असे समजा. प्रत्येक~$n\geq1$ साठी
  जगांचा निश्चित संच $W=\{1,\ldots,n\}$ आणि वैश्विक संबंध
  $R=W\times W$ घेऊ. $p_1$, \dots, $p_k$ यांच्या $2^{nk}$
  मूल्यांकनांची यादी सांत आहे; इतर सर्व विधानचलांना रिक्त संच नेमू.
  त्यामुळे प्रत्येक आकारासाठी तपासायच्या वैश्विक प्रतिमानांची सांत
  प्रभावी यादी मिळते.
  म्हणून दोन प्रक्रिया \emph{समांतर} चालवू शकतो:
  एका बाजूने सिद्धता पद्धतशीरपणे तयार करून
  \Log{S5} ची सर्व प्रमेये क्रमाने नोंदवायची;
  दुसऱ्या बाजूने $1$, $2$, \dots जगे असलेली सर्व
  वैश्विक प्रतिमाने क्रमाने नोंदवायची आणि त्यांपैकी
  एखाद्या प्रतिमानातील एखाद्या जगावर~$!A$ असत्य
  आहे का हे तपासायचे. अखेरीस या दोनपैकी एक
  प्रक्रिया उत्तर देईल. कारण
  \olref[com][fra]{thm:generaldet} आणि
  \olref[fmp]{cor:S5fmp} नुसार, एकतर~$!A$
  !!{derivable} आहे किंवा ते एखाद्या सांत
  वैश्विक प्रतिमानात असत्य आहे.
\end{proof}

वरील सिद्धता \Log{S5} साठी चालते, कारण
वैश्विक प्रतिमानांची निस्यंदने आपोआप वैश्विक
असतात. परावर्तित्व आणि सर्वत्र उत्तरवर्तित्व
यांबाबतही तसेच आहे; इतर गुणधर्मांसाठी अधिक
काम करावे लागते.

\end{document}
